\chapter{CONCLUSION AND FUTURE ENHANCEMENTS}

\section{Conclusion}

This case study compared six multiclass classifiers on the UCI Dry Bean dataset
under a single stratified 80/20 split, cross-validated hyperparameter search on
the training partition and one held-out evaluation. The \BestModel{} was selected
on cross-validated macro F1 (\BestCV{} $\pm$ \BestCVStd{}) and reached \BestAcc{}
accuracy and \BestFOne{} macro F1 on the \TestRows{} held-out beans, with a
one-vs-rest ROC-AUC of \BestAUC.

The ranking is informative rather than surprising. The five models with nonlinear
or locally adaptive boundaries all land within 0.007 macro F1 of each other, and
the linear baseline sits inside that band rather than below it. Once a learner can
express a curved boundary, the remaining differences come from how efficiently it
uses 10,888 training beans, not from the shape of the hypothesis class. The narrow
spread also means the winner should not be over-interpreted: a different split
seed could reorder the middle of the table.

All seven objectives were met. The morphological analysis identified size as the
separator for Bombay and shape ratios as the separator inside the remaining cloud.
Validation found a complete, imbalanced dataset with 21 strongly correlated
feature pairs, and drove the preprocessing decisions. The preprocessing and split
were implemented as pipelines refitted inside every fold, and the test suite
asserts that the scaler is estimated from training data only. Six classifiers were
tuned within a shared protocol and compared on accuracy, macro precision, recall
and F1, weighted F1, ROC-AUC, confusion matrices, per-class scores and cost.
Selection used cross-validated macro F1 with the held-out result, per-class
behaviour and cost as corroboration, and the demonstration is a saved pipeline
with a command-line interface that reproduces the reported predictions.

Three findings are worth carrying forward. Feature selection did not help:
dropping five of the 16 columns changed cross-validated macro F1 from
\AblationFull{} to \AblationReduced{}, inside the fold-to-fold noise. Balancing the
class weights did not help either (\CWNone{} against \CWBalanced{}), which
contradicts the assumption behind much of the published resampling work on this
dataset: the minority class is small but already well separated. And the errors
are concentrated: \TestErrors{} of \TestRows{} beans were misclassified and
\TopConfusion{} accounts for \TopConfusionN{} of them, so the practical limitation
of this model is one confusable pair rather than a diffuse weakness.

\section{Limitations}

\textbf{Single dataset and single split.} All beans were imaged by one laboratory
with one camera and one segmentation procedure, and the evaluation uses one
stratified split with a fixed seed. The reported accuracy is therefore an internal
estimate: a second laboratory would require recalibration, and repeated or nested
cross-validation would give a tighter confidence interval than the single
cross-validated spread in Table~\ref{tab:cv_results}.

\textbf{Retained duplicates and outliers.} \DuplicateRows{} rows repeat an earlier
row and were kept because the features are quantised pixel counts; large beans were
kept because they are genuine Bombay specimens. Both decisions are defensible but
neither is neutral: the effective sample size is slightly below \DatasetRows{}, and
the model's behaviour on an unusually small Bombay is untested.

\textbf{Class overlap is irreducible with these features.} Cali, Barbunya, Sira and
Dermason overlap in the 16-dimensional space, and the remaining errors are
geometric rather than algorithmic. Colour, texture or three-dimensional shape would
be needed to separate them reliably.

\textbf{Importance measures are not causal.} The permutation and SHAP rankings
describe how this fitted model uses its inputs. Because 21 feature pairs correlate
above 0.90, credit is shared arbitrarily between near-duplicate columns, and a high
rank for \TopPermFeature{} does not establish that the ratio causes the variety
label.

\textbf{Scope of the comparison.} Six classical learners were compared. Deep
models, stacking, calibrated ensembles and cost-sensitive variants were not
exhausted, so the conclusion is a statement about the six families tested.

\section{Future Enhancements}

Repeating the split with several seeds, or nesting the hyperparameter search
inside an outer cross-validation loop, would separate split variance from model
variance and give a defensible interval on the reported accuracy. Obtaining images
from a second laboratory would test whether the features survive a change of camera
and segmentation.

The most direct modelling improvement is to add colour and texture descriptors.
The current features are purely geometric, and the confusable group differs visibly
in colour even where shape overlaps. A two-stage model, first separating the
size-separated varieties and then resolving the overlapping group, would match the
structure revealed by the EDA. Calibrating the Platt-scaled probabilities with
isotonic regression would make the confidence values usable for a rejection option,
in which uncertain beans are referred to a human inspector.

Finally, the saved pipeline is small enough to run on a single-board computer
attached to a grading line. A latency and memory budget, together with a drift
monitor on the input distribution, would be the next engineering step.
